\section{分析与局限}
\label{sec:analysis}

核心主张是校准存储与读取之间的关系。来源片段是被保存的观测；身份族、关系或版本是推断得到的访问结构；检索到的段落则是智能体主动选择检查的内容。把这些类别分开可以支持修订，但不能保证身份正确，也不能保证读到的段落蕴含答案。来源门只验证来源，不验证身份或蕴含关系。

这种存储选择在提出具体问题之前就决定了信息能否保留。只保存摘要会提前压缩历史，后来的问题可能需要摘要中被省略的限定条件、旧状态或原始措辞。Deep-Dream保留来源观测，把图结构只用于定位，因此当前值仍然可以回到产生它的历史观测，而不是覆盖这些观测。

延迟对齐让这种表示可以修订。过早合并可能把两个同名的人放在同一个地址下，污染双方的历史；暂时拆分只增加候选，之后可以通过重定向链接来修正。因此，入库实验说明的是成本和可修订性变化，并不能证明当前身份标签已经完全正确。版本链接还保留了当前值的形成过程：处理时间排列存储状态，原文中提到的日期仍属于观测内容。查询可以要求最新值，也可以要求较早的值，而不必改写历史记录。

Align+使写入调用减少46--65\%，但其中同时包含延迟对齐、批处理和并行处理。同名族在MemoryAgentBench上由340降至124，在LongMemEval上由482降至99。同名重复代理不是身份标签，配对质量区间为$-7$至$+13$个百分点。MCC和FC-MH各提高3个百分点，95\%区间为$-7$至$+13$，因此这些结果说明的是构建成本与可修订性的权衡，而不是身份准确率提升。熵论证还以固定目标和保留前驱状态为条件。

MAB路径比较中，固定检索、记忆工具和原文阅读在1,074道题上的Kimi-K3得分分别为33.3\%、40.0\%和71.0\%。三条完整路径在工具、原文访问、迭代方式、证据处理和预算上均有差异，因此不能据此单独分离身份层或图结构的作用。原文阅读在相同的614题上为70.8\%，全文上下文为78.7\%；另有460道题超出全文预算。离线图消融没有稳定收益。即使活动上下文较小，反复检索也可能提高成本；来源检查也不能解决停止时机或答案支持问题。

MEME在694次后任务检查中的100个删除检查上准确率为0\%，级联和缺失准确率分别为46.34\%和59.23\%。该路径没有目标级墓碑标记和传播式抑制。LongMemEval-S在Qwen3.7-plus同时担任回答者和评审时为84.6\%，协议与其他结果不同。LoCoMo的95.19\%和94.22\%是历史单次诊断，其中1,540行有1,529行早于运行时代码指纹记录。这些结果要求直接标注身份、检查蕴含关系，并开展资源匹配的读取实验。

\paragraph{尚未测量的内容。} 通过来源检查的段落仍可能与问题无关，正确链接的身份也可能需要另一份文档。当前实验没有带标签的身份决策、蕴含标注或资源匹配的停止曲线，删除结果也没有测试完整的抑制设计。后续评估应标注候选合并，在读取策略之间固定调用数和token数，并分别评价当前查询与历史查询。
